% 这是中国科学院大学计算机科学与技术专业《计算机组成原理（研讨课）》使用的实验报告 Latex 模板
% 本模板与 2024 年 2 月 Jun-xiong Ji 完成, 更改自由 Shing-Ho Lin 和 Jun-Xiong Ji 于 2022 年 9 月共同完成的基础物理实验模板
% 如有任何问题, 请联系: jijunxoing21@mails.ucas.ac.cn
% This is the LaTeX template for report of Experiment of Computer Organization and Design courses, based on its provided Word template. 
% This template is completed on Febrary 2024, based on the joint collabration of Shing-Ho Lin and Junxiong Ji in September 2022. 
% Adding numerous pictures and equations leads to unsatisfying experience in Word. Therefore LaTeX is better. 
% Feel free to contact me via: jijunxoing21@mails.ucas.ac.cn

\documentclass[11pt]{article}

\usepackage[a4paper]{geometry}
\geometry{left=2.0cm,right=2.0cm,top=2.5cm,bottom=2.5cm}

\usepackage{ctex} % 支持中文的LaTeX宏包
\usepackage{amsmath,amsfonts,graphicx,subfigure,amssymb,bm,amsthm,mathrsfs,mathtools,breqn} % 数学公式和符号的宏包集合
\usepackage{algorithm,algorithmicx} % 算法和伪代码
\usepackage[noend]{algpseudocode} % 算法和伪代码
\usepackage{fancyhdr} % 自定义页眉页脚
\usepackage[framemethod=TikZ]{mdframed} % 创建带边框的框架
\usepackage{fontspec} % 字体设置
\usepackage{adjustbox} % 调整盒子大小
\usepackage{fontsize} % 设置字体大小
\usepackage{tikz,xcolor} % 绘制图形和使用颜色
\usepackage{multicol} % 多栏排版
\usepackage{multirow} % 表格中合并单元格
\usepackage{pdfpages} % 插入PDF文件
\usepackage{listings} % 在文档中插入源代码
\usepackage{wrapfig} % 文字绕排图片
\usepackage{bigstrut,multirow,rotating} % 支持在表格中使用特殊命令
\usepackage{booktabs} % 创建美观的表格
\usepackage{circuitikz} % 绘制电路图
\usepackage{zhnumber} % 中文序号（用于标题）
\usepackage{tabularx} % 表格折行

\definecolor{dkgreen}{rgb}{0,0.6,0}
\definecolor{gray}{rgb}{0.5,0.5,0.5}
\definecolor{mauve}{rgb}{0.58,0,0.82}
\lstset{
  frame=tb,
  aboveskip=3mm,
  belowskip=3mm,
  showstringspaces=false,
  columns=flexible,
  framerule=1pt,
  rulecolor=\color{gray!35},
  backgroundcolor=\color{gray!5},
  basicstyle={\small\ttfamily},
  numbers=none,
  numberstyle=\tiny\color{gray},
  keywordstyle=\color{blue},
  commentstyle=\color{dkgreen},
  stringstyle=\color{mauve},
  breaklines=true,
  breakatwhitespace=true,
  tabsize=3,
}

% 超链接支持 (一定要在大多数包之后, cleveref 之前)
% colorlinks=true 使链接文字着色而不是加方框; hidelinks 可去掉颜色
\usepackage[colorlinks=true,linkcolor=blue,citecolor=blue,urlcolor=magenta]{hyperref}

% 轻松引用, 可以用\cref{}指令直接引用, 自动加前缀. 需要在 hyperref 之后加载
% 例: 图片label为fig:1  => \cref{fig:1} -> Figure.1; \ref{fig:1} -> 1 (纯数字)
\usepackage[capitalize]{cleveref}
% \crefname{section}{Sec.}{Secs.}
\Crefname{section}{Section}{Sections}
\Crefname{table}{Table}{Tables}
\crefname{table}{Table.}{Tabs.}

% \setmainfont{Palatino Linotype.ttf}
% \setCJKmainfont{SimHei.ttf}
% \setCJKsansfont{Songti.ttf}
% \setCJKmonofont{SimSun.ttf}
\punctstyle{kaiming}
% 偏好的几个字体, 可以根据需要自行加入字体ttf文件并调用

\renewcommand{\emph}[1]{\begin{kaishu}#1\end{kaishu}}

% 对 section 等环境的序号使用中文
\renewcommand \thesection{\zhnum{section}、}
\renewcommand \thesubsection{\arabic{section}.\arabic{subsection}}


%%%%%%%%%%%%%%%%%%%%%%%%%%%
%改这里可以修改实验报告表头的信息
\newcommand{\name}{寇逸欣}
\newcommand{\studentNum}{2023K8009922004}
\newcommand{\major}{计算机科学与技术}
\newcommand{\labNum}{01}
\newcommand{\labName}{信息熵计算与齐夫定律的验证}
%%%%%%%%%%%%%%%%%%%%%%%%%%%

\begin{document}

\input{tex_file/head.tex}

\section{实验目的}
本实验旨在通过实际编程操作，加深对自然语言统计特性的理解。具体目标包括：
\begin{enumerate}
    \item 学习网络爬虫技术，从互联网收集中英文双语语料。
    \item 对收集的文本数据进行清洗和预处理。
    \item 计算不同规模下英语（字母、单词）和汉语（汉字、词汇）的概率分布及信息熵。
    \item 验证齐夫定律（Zipf's law）。
    \item 探究样本规模的变化对统计结果（概率、信息熵）的影响。
\end{enumerate}

\section{数据收集与预处理}
\subsection{数据爬取设计}
本次实验使用了 \textbf{Scrapy} 爬虫框架。作为数据来源：
\begin{itemize}
    \item \textbf{英文语料}：主要从 BBC 等英文新闻网站爬取新闻报道文本。
    \item \textbf{中文语料}：主要从新浪新闻等中文网站爬取新闻报道文本。
\end{itemize}
通过使用 Scrapy 开源框架，我们提取了新闻标题和正文部分，并将结果存储为 JSON 格式。

\subsection{数据清洗}
本实验针对中英文语料的结构特征和噪声情况，在 Scrapy 的 Pipeline 中使用正则表达式和字符转义等技术设计了专门的清洗策略：
\begin{enumerate}
    \item \textbf{基础与结构化噪音过滤：}
    \begin{itemize}
        \item HTML 实体转义：利用 \texttt{html.unescape} 技术，将网页文本中的 \verb|&nbsp;| 等转义字符还原。
        \item 移除网页标签与超链接：通过正则表达式严格移除非文本结构的 HTML 标签，并滤除形如 \texttt{http(s)://...} 的 URL 链接。
    \end{itemize}
    \item \textbf{中文文本处理流程：}
    \begin{itemize}
        \item 字符集白名单过滤：通过正则表达式限制允许出现的字符区间，仅保留基础汉字（\verb|\u4e00-\u9fff| 等）、字母数字，以及部分中常用中文标点符号（如破折号、书名号等）。
        \item 非法字符拦截：所有未包含在白名单中的乱码或特殊表情包字符均替换为空格。后续再传递给 jieba 进行分词截断。
    \end{itemize}
    \item \textbf{英文文本处理流程：}
    \begin{itemize}
        \item 专用字符白名单约束：仅保留英文字母、数字及基础英文标点等，清理排版时的冗余符号。
        \item 大小写与空白符标准化：采用 \verb|\s+| 统一压缩包括制表符、连续回车在内的大量空白字符为单一空格，并将内容规范化为小写，确保独立的单词计数的纯净度。
    \end{itemize}
    \item \textbf{短文本丢弃检查：}
    经过格式化清洗后，在最终保留阶段引入基于长度的 \texttt{DropItem} 机制。一旦清洗后正文总长度不足 30 个字符，或提取不到有效正文和标题的记录，将被判定为无效资讯或非结构文本并直接丢弃，从根源保证语料的总体纯度。
\end{enumerate}

\subsection{语料库规模}
为了进行不同维度的对比分析，我们将清洗后的数据按照文章数量划分成了不同的规模类别。各规模下的中英文语料统计信息如表~\ref{tab:corpus_size} 所示：

\begin{table}[htbp]
    \centering
    \caption{不同规模下的英中文语料统计信息}
    \resizebox{\textwidth}{!}{
    \begin{tabular}{c|cccc|cccc}
        \toprule
        \multirow{2}{*}{\textbf{目标规模}} & \multicolumn{4}{c|}{\textbf{英文语料}} & \multicolumn{4}{c}{\textbf{中文语料}} \\
        \cmidrule{2-9}
         & \textbf{文章数} & \textbf{总字母数} & \textbf{总词数(Tokens)} & \textbf{词库大小(Types)} & \textbf{文章数} & \textbf{总汉字数} & \textbf{总词数(Tokens)} & \textbf{词库大小(Types)} \\
        \midrule
        100 & 121 & 414999 & 87711 & 8597 & 126 & 105092 & 55669 & 11475 \\
        200 & 220 & 843835 & 178694 & 12908 & 230 & 218472 & 116306 & 17775 \\
        300 & 318 & 1092917 & 231313 & 13908 & 328 & 337196 & 179473 & 23204 \\
        500 & 521 & 1596067 & 347487 & 17344 & 529 & 573263 & 308032 & 31342 \\
        1000 & 1011 & 3396189 & 721787 & 27095 & 1027 & 1250475 & 664531 & 50772 \\
        2000 & 2020 & 5972797 & 1286905 & 36796 & 2029 & 2212022 & 1187339 & 65265 \\
        \bottomrule
    \end{tabular}
    }
    \label{tab:corpus_size}
\end{table}

\section{信息熵计算与对比分析}
信息熵是衡量信息量大小和系统不确定性的重要指标，其计算公式为：
$$ H(X) = - \sum_{x \in X} P(x) \log_2 P(x) $$
其中，$X$ 为离散随机变量的取值空间，$P(x)$ 为对应事件发生的概率。

在本次实验中，我们分别统计了中英文语料在不同样本规模下的语言单元（英文字母、英文单词、汉字、中文词汇）的频数，从而估计各自的概率分布 $P(x)$，并计算出相应的信息熵。具体计算结果如表~\ref{tab:entropy} 所示。

\begin{table}[htbp]
    \centering
    \caption{不同样本规模下的中英文信息熵计算结果 (bits/symbol)}
    \begin{tabular}{ccccc}
        \toprule
        \textbf{样本规模(文章数)} & \textbf{英文字母} & \textbf{英文单词} & \textbf{汉字} & \textbf{中文词汇} \\
        \midrule
        100 & 4.1715 & 9.9429 & 9.5685 & 11.4394 \\
        200 & 4.1708 & 10.1239 & 9.5838 & 11.6135 \\
        300 & 4.1678 & 10.0770 & 9.5670 & 11.6797 \\
        500 & 4.1813 & 10.1675 & 9.5843 & 11.7409 \\
        1000 & 4.1795 & 10.3970 & 9.6256 & 12.0164 \\
        2000 & 4.1806 & 10.4322 & 9.6447 & 12.0063 \\
        \bottomrule
    \end{tabular}
    \label{tab:entropy}
\end{table}

\subsection{结果分析与讨论}
结合表~\ref{tab:entropy} 中的数据，我们可以分级、分阶段得出以下不同维度的横纵向规律与结论：

\paragraph{一、 字符级别的信息熵（英文字母与汉字）呈现高度稳定性}
\begin{enumerate}
    \item \textbf{数据表现：} 英文字母的信息熵基本稳定在 4.17 至 4.18 bits/symbol 左右，汉字的信息熵也稳定在 9.56 至 9.64 bits/symbol 之间，无论在 100 篇还是 2000 篇语料时都没有量级上的剧烈跳动。
    \item \textbf{理论解释：} 字母和汉字作为语言的基本书写符号，有着本质区别（英文 26 字母为封闭集，汉字为多达数千的半开放集）。但因为常用字在达到一定体量的自然语料中，其使用频率和分布很快趋于极限的帕累托稳定，所以很难因文本数成倍放大而带入大量未见字符进行分布重塑。
    \item \textbf{交叉验证：} 从实验计算所得的汉字平均熵约为 9.6 bits 可见，这与中国国家标准 GB2312 编码的汉字分布以及主流自然语言学的理论认知吻合。
\end{enumerate}

\paragraph{二、 词汇级别的信息熵总体呈现增长趋势，但在大规模时出现局部振荡}

从表~\ref{tab:entropy} 的中文词汇熵（11.43 $\to$ 11.61 $\to$ 11.67 $\to$ 11.74 $\to$ 12.01 $\to$ 12.00）及英文单词熵（9.94 $\to$ 10.12 $\to$ 10.07 $\to$ 10.16 $\to$ 10.39 $\to$ 10.43）中，我们可以明显分为两个典型的阶段现象来分析：
\begin{enumerate}
    \item \textbf{典型增长阶段（例如从 100 篇到 200 篇）：}
    \begin{itemize}
        \item \textbf{数据表现：} 语料增加带来了独立词数剧增，以英文 500 到 1000 篇为例，独立单词从 17344 上涨至 27095。
        \item \textbf{分布变化：} 这是符合预期的增长。文本规模扩大，且由于是从新浪或 BBC 不断刷新、包含不同板块主题的抓取，引入了更多专业级和未见过的命名实体。关键证据是，长尾效应变得明显，稀释了“the/of”或“的/了/在”这类高频词原有的统治级概率。
        \item \textbf{结论：} 稀疏新词的扩充导致最常见词的集中度被打破，分布更倾向于“均匀长尾”。从而增大了该体系的统计不确定性，信息熵因此水涨船高。
    \end{itemize}
    \item \textbf{局部振荡与收敛阶段（例如中文 1000 篇到 2000 篇）：}
    \begin{itemize}
        \item \textbf{数据表现：} 中文 1000 到 2000 篇时信息熵由 12.01 微跌至 12.00 bits，呈现出轻微的上下振荡波动。
        \item \textbf{理论解释：} 根据公式 $H = -\sum P(x) \log_2 P(x)$，高频词对整体熵的数值牵引占主导地位。由于长尾极低频新词所贡献的微小熵增已相当有限，如果抽样导致占比极大的头部高频词出现极其微小的上扬时，就会在数学上抵消甚至压倒字典扩充带来的熵增。因此，熵值在末期微跌和横跳并不异常，这代表了在大规模语料库中，词频分布的稳定性开始显现，信息熵趋于收敛的迹象逐渐浮现。
    \end{itemize}
\end{enumerate}

\section{齐夫定律 (Zipf's Law) 的验证}
齐夫定律指出，在自然语言预料中，一个单词出现的频率 $f$ 与其在频率表中的排名 $r$ 成反比，即：
$$ f \times r \approx \text{常数} $$
或者在对数坐标系下表现为近似一条斜率为 -1 的直线：
$$ \log(f) = \log(C) - \log(r) $$

利用我们收集到的多规模英文文本数据，对其进行词频统计并绘制了 $\log(r) - \log(f)$ 对数图。

\begin{figure}[htbp]
    \centering
    \subfigure[规模 100]{\includegraphics[width=0.32\textwidth]{zipf_law_en_100.png}}
    \subfigure[规模 200]{\includegraphics[width=0.32\textwidth]{zipf_law_en_200.png}}
    \subfigure[规模 300]{\includegraphics[width=0.32\textwidth]{zipf_law_en_300.png}}
    \\
    \subfigure[规模 500]{\includegraphics[width=0.32\textwidth]{zipf_law_en_500.png}}
    \subfigure[规模 1000]{\includegraphics[width=0.32\textwidth]{zipf_law_en_1000.png}}
    \subfigure[规模 2000]{\includegraphics[width=0.32\textwidth]{zipf_law_en_2000.png}}
    \caption{不同样本规模下英文单词的 Zipf 定律验证对数图}
    \label{fig:zipf_law}
\end{figure}

如图~\ref{fig:zipf_law} 所示，从对数坐标的散点图及拟合结果可以看出，所有规模语料均表现出词频与排名近似线性相关的特性。结合图表，我们得出以下实验结论与偏差分析：
\begin{enumerate}
    \item \textbf{基本符合理论：} 随着数据规模的增加，主体部分的线性拟合越发平滑。所有样本的分布都很好地契合了 Zipf's law 标准斜率近似为 -1 的参考线，表明词频与排名的对数关系具有较高的线性相关性。
    \item \textbf{高频词区域（首部偏差）：} 在所有图中，曲线的头部（最左侧高频词区域）都存在一定程度高于或偏离期望直线的现象。这主要是因为大量的冠词、介词等功能词（如 "the", "of", "a"）占据了频率的顶端，其出现频率远超普通内容词，导致了分布前段的陡峭化。
    \item \textbf{低频词区域（尾部阶梯化）：} 曲线的尾部（最右侧）出现了明显的向下弯曲和水平的“台阶（Step）”状结构。这种现象源于以下几个原因：
    \begin{itemize}
        \item \textbf{有限样本稀疏性：} 在有限的自然语料中，大量稀有词或生僻词仅出现了一次或极少数次，导致在对数坐标尾部的相同频率点上聚集了大量不同排名的词，从而形成离散的水平“台阶”。
        \item \textbf{数据离散化放大：} 排名越靠后，词频绝对数值越低，频率的下降粒度只能是整数步进（不再是连续分布）。这种低维度的整数截断跳变效应在对数坐标图上被视觉尺度放大了。
    \end{itemize}
\end{enumerate}

\section{结论与心得}
本实验完整地执行了从网络爬虫收集语料、数据清洗、特征统计到可视化与理论验证的全流程。
\begin{itemize}
    \item 理论上，直观地认识到了信息熵作为衡量语言信息量大小的有效指标。并成功在英语语料上验证了语言学的重要假设——齐夫定律。
    \item 工程实现上，熟练度掌握了 Scrapy 框架的使用以及文本数据预处理的方法。
\end{itemize}

\end{document}
